\chapter{Matrix-free action and Lanczos solving}
\label{chap:lanczos}

\section{Apply the same terms without storing them}

Iterative eigensolvers only require products $y=Hx$.  For a vector
$x_j=\braket{s_j}{\psi}$, the matrix-free loop is the sparse assembly loop with
one change: instead of appending $H_{ij}$, accumulate
\begin{equation}
y_i \mathrel{+}= H_{ij}x_j.
\end{equation}
The basis and state-index dictionary are retained, but the CSR arrays are not.
Correctness is checked by applying both implementations to the same random
real and complex vectors.

\lstinputlisting[caption={CSR versus matrix-free multiplication.},label={lst:matrixfree}]{code/step04_matrix_free.py}

Using the same elementary hopping routine in both paths is intentional.  It
reduces the risk that an optimized implementation and a reference
implementation silently adopt different sign conventions.

\section{Lanczos through a linear operator}

SciPy's \code{LinearOperator} wraps the matrix-free action and supplies it to
\code{eigsh}.  For a Hermitian Hamiltonian, requesting
\code{which="SA"} returns the algebraically smallest eigenvalues.  A
deterministic random initial vector is preferable to a constant vector because
a constant vector may lie entirely within one symmetry block and miss a lower
state in another block.

ARPACK requires the number of requested eigenvalues $k$ to satisfy
$k<\basisdim$.  Tiny sectors and complete-spectrum requests below a strict
size threshold therefore use dense diagonalization.  Dense fallback is an
edge-case convenience, not a route for production sectors.

\section{Residuals are the final contract}

An eigensolver's convergence flag is not enough.  Normalize every returned
vector and recompute
\begin{equation}
r_a=\norm{H\psi_a-E_a\psi_a}_2
\label{eq:residual}
\end{equation}
with the matrix-free action.  The default ARPACK tolerance is $10^{-10}$;
typical absolute residuals in the tutorial examples are between $10^{-15}$
and $10^{-10}$.

\lstinputlisting[caption={Low-energy states and independent residual checks.},label={lst:groundstate}]{code/step05_ground_state.py}

Degenerate eigenvalues require care.  Individual eigenvectors within a
degenerate subspace are not unique, so compare the eigenvalue multiset or the
projector onto the subspace rather than component-by-component vectors.

\paragraph{Checkpoint.}
Test CSR and matrix-free products on seeded random real and complex vectors,
compare their low eigenvalues, and require a small independently recomputed
residual for every eigenpair used in a numerical assertion.
